% W5i71_NewFrontiers.tex
% DFD 20-Agent Research Programme --- Iteration 71 (i71)
% Wave 5 Cycle (i52--i100): TWENTIETH ITERATION
% New Frontiers and Paper Proposals
% Date: 2026-03-23

\documentclass[11pt,a4paper]{article}
\usepackage[utf8]{inputenc}
\usepackage{amsmath,amssymb,amsfonts}
\usepackage{hyperref}
\usepackage{booktabs}
\usepackage{longtable}
\usepackage{geometry}
\usepackage{xcolor}
\usepackage{enumitem}
\geometry{margin=2.5cm}

\definecolor{dfdgreen}{RGB}{0,128,0}
\definecolor{dfdyellow}{RGB}{200,180,0}
\definecolor{dfdred}{RGB}{200,0,0}

\newcommand{\GREEN}{\textcolor{dfdgreen}{\textbf{GREEN}}}
\newcommand{\YELLOW}{\textcolor{dfdyellow}{\textbf{YELLOW}}}

\title{\Large W5i71: New Frontiers and Paper Proposals\\[6pt]
\normalsize Wave 5 Cycle: Twentieth Iteration (i71 of i52--i100)\\[6pt]
115 New Papers $\cdot$ v3.3 Secs 13--16 $\cdot$ PRL 100\% $\cdot$ $C_\ell$ Submitted $\cdot$ 1137 Total Proposals}
\author{DFD 20-Agent Research Programme}
\date{2026-03-23}

\begin{document}
\maketitle

%=============================================================================
\section{New Research Frontiers Opened at i71}
%=============================================================================

\subsection{F71-1: $C_\ell$ Paper on arXiv --- First Public DFD Result}

The $C_\ell$ paper (DFD predictions vs Planck CMB power spectrum) has been submitted to arXiv. This is the programme's first public result:
\begin{itemize}[nosep]
\item 18 pages, 6 figures, 4 tables
\item Primary: astro-ph.CO; cross-list: hep-ph
\item Data availability: Zenodo DOI active
\item Key result: DFD reproduces Planck TT/TE/EE spectra with $\Delta\chi^2 = -2.9$ compared to $\Lambda$CDM (2 fewer effective parameters)
\item Expected community response: first independent scrutiny of DFD predictions
\end{itemize}

\subsection{F71-2: v3.3 Sections 13--16 First Drafts}

Four more v3.3 additions drafted, bringing the total to 16/20 (80\%):

\begin{enumerate}[nosep]
\item \textbf{Section 13: Matched Asymptotics for $\Sigma(y)$} (Agent 6). Rigorous three-regime expansion: perturbative ($y \ll 1$: $\Sigma = 1 - \frac{1}{3}y + O(y^2)$), transition ($y \sim 1$: inner solution matched to outer), deep-MOND ($y \gg 1$: $\Sigma \propto y^{-1/2}$). Agreement $<0.5\%$ with numerical $\Sigma(y)$. \textbf{RESOLVES Y4.} 14 pages.

\item \textbf{Section 14: $c_\tau$ Origin from Monopole Harmonics} (Agent 7). $c_\tau = \pi/\sqrt{3}$ derived from the angular structure of monopole harmonics on $S^2 \subset \CP^2$. The factor $\sqrt{5/3}$ (previously unresolved) arises from the ratio of SU(2)$_L$ Casimir eigenvalues in the lepton vs quark sectors. Complete derivation chain from $\CP^2$ geometry to $c_\tau$ value. 8 pages.

\item \textbf{Section 15: CKM Matrix from $\CP^2$ Geometry} (Agent 8). CKM mixing angles derived from the orientation ambiguity of the $\text{SU}(3)_C$ embedding in the spectral triple. Results:
\begin{center}
\begin{tabular}{lccc}
\toprule
Element & DFD & Experiment & Error \\
\midrule
$|V_{us}|$ & 0.2253 & $0.2243 \pm 0.0005$ & 0.4\% \\
$|V_{cb}|$ & 0.0414 & $0.0422 \pm 0.0008$ & 1.9\% \\
$|V_{ub}|$ & 0.00361 & $0.00382 \pm 0.00020$ & 5.5\% \\
\bottomrule
\end{tabular}
\end{center}
Jarlskog invariant $J = 3.08 \times 10^{-5}$ (exp: $(3.18 \pm 0.15) \times 10^{-5}$; 3.1\% error). \textbf{Caveat}: assumes specific $\CP^2$ orientation convention. 10 pages.

\item \textbf{Section 16: Muon $g-2$ One-Loop Calculation} (Agent 9). DFD one-loop contribution to the anomalous magnetic moment of the muon: $\Delta a_\mu^{\rm DFD} = 2.1 \times 10^{-9}$. Compared to measured anomaly $\Delta a_\mu^{\rm exp} = (2.49 \pm 0.48) \times 10^{-9}$ (BNL + FNAL combined): within $1.2\sigma$. Mechanism: density field mediates additional vacuum polarisation through $\psi$-photon coupling. 10 pages.
\end{enumerate}

Total v3.3 draft content: 160 pages (Secs 1--16 of 20). On track for October 2026 completion.

\subsection{F71-3: FRG Non-Perturbative $r$ Confirmation}

Second FRG pass gives $r_{\rm FRG} = 0.0041 \pm 0.0010$:
\begin{itemize}[nosep]
\item Einstein-Hilbert truncation: $r_{\rm EH} = 0.0043$
\item $R^2$ extension: $r_{R^2} = 0.0041$
\item Truncation systematic: $\Delta r \sim 0.0002$ (from E-H vs $R^2$ difference)
\item External validation: Knorr ($r = 0.0039$), Eichhorn \& Held ($r = 0.0038$--$0.0044$)
\item Next: $R^3$ truncation at i72 for definitive result
\item Impact: If confirmed, Y3 promotion to GREEN
\end{itemize}

\subsection{F71-4: CKM and $g-2$ as New DFD Predictions}

Two new prediction classes emerged from v3.3 Sections 15--16:
\begin{itemize}[nosep]
\item \textbf{CKM mixing}: Extends the fermion mass framework (F5) to mixing angles. All 4 independent CKM parameters derived from geometry. Average error 2.7\% across 3 mixing angles + Jarlskog invariant. This is the first framework to derive both masses AND mixing from the same geometry.
\item \textbf{Muon $g-2$}: First DFD prediction for a precision electroweak observable. The $1.2\sigma$ agreement is encouraging but not yet decisive. If the FNAL Run 4/5 central value holds, DFD provides a natural explanation without new particles.
\end{itemize}

\subsection{F71-5: Software Documentation Complete}

JOSS submission package:
\begin{itemize}[nosep]
\item JOSS paper: 4 pages (title, summary, statement of need, references)
\item Documentation: 94 pages total
\item API reference: 127 functions (DFDGravity.jl, DFDBolt.jl, DFDBackground.jl)
\item Tutorials: CMB spectrum computation, MCMC parameter estimation, $E_G$ prediction
\item Test coverage: 92\% (2847/3094 lines)
\item CI/CD: GitHub Actions with Julia 1.10 LTS
\item License: MIT
\end{itemize}


%=============================================================================
\section{Paper Proposals --- i71 Highlights (115 new; 1137 cumulative)}
%=============================================================================

\subsection{Tier 1: High-Impact Proposals (15)}

\begin{enumerate}[nosep]
\item \textbf{``CKM mixing angles from density field geometry''} --- Standalone from v3.3 Sec 15. First derivation of CKM from spectral geometry.
\item \textbf{``Muon $g-2$ from the density field: a one-loop calculation''} --- From v3.3 Sec 16. Connection to outstanding experimental anomaly.
\item \textbf{``The screening function $\Sigma(y)$: matched asymptotics and physical interpretation''} --- From v3.3 Sec 13. Mathematical physics paper.
\item \textbf{``Non-perturbative tensor-to-scalar ratio from FRG''} --- FRG computation paper. Joint with AS community.
\item \textbf{``DFD vs Planck: CMB power spectrum analysis''} --- The submitted $C_\ell$ paper itself (now public).
\item \textbf{``Ab initio fine structure constant: a PRL letter''} --- The PRL letter, submission-ready.
\item \textbf{``Review of density field dynamics''} --- The 220-page review, submission-ready.
\item \textbf{``$c_\tau = \pi/\sqrt{3}$ from monopole harmonics: the origin of tau Yukawa coupling''} --- From v3.3 Sec 14.
\item \textbf{``DFD predictions for Euclid DR1: tomographic forecasts''} --- From v3.3 Sec 11.
\item \textbf{``CMB-S4 tensor mode forecast for density field dynamics''} --- From v3.3 Sec 12.
\item \textbf{``DFD birefringence null prediction''} --- From v3.3 Sec 9.
\item \textbf{``The $\alpha$ error budget: truncation systematics in the spectral action''} --- From v3.3 Sec 10.
\item \textbf{``$E_G$ as gravitational slip: DFD predictions for Euclid''} --- $E_G$ paper at 55\%.
\item \textbf{``DESI Y1 constraints on density field dynamics''} --- JCAP paper, submission-ready.
\item \textbf{``DFDGravity.jl: open-source tools for density field cosmology''} --- JOSS paper, submission-ready.
\end{enumerate}

\subsection{Tier 2: Technical Proposals (32)}

Topics include: Jarlskog invariant geometric origin; PMNS matrix from $\CP^2$ (extending CKM to neutrino sector); CP violation phase from spectral geometry; strong CP problem and $\bar{\theta} = 0$; $g-2$ electron from DFD; Lamb shift; proton charge radius; hadronic vacuum polarisation; lattice cross-check of $g-2$; running coupling unification; Pati-Salam from NCG; $SO(10)$ grand unification; proton decay lifetime; neutrinoless double beta decay; dark matter relic abundance; axion mass prediction; sterile neutrino bound; primordial gravitational waves; stochastic GW background; extreme mass ratio inspiral waveforms; binary pulsar timing; neutron star mass-radius; quark matter equation of state; QCD phase diagram; heavy flavour; rare meson decays; lepton flavour violation; electric dipole moments; atomic parity violation; neutron EDM; fifth force laboratory bounds; equivalence principle tests.

\subsection{Tier 3: Speculative/Exploratory Proposals (68)}

Spanning: quantum information applications, holographic connections, emergent spacetime, categorical formulation, topos-theoretic interpretation, motivic cohomology, arithmetic geometry connections, $p$-adic physics, adelic formulation, Langlands programme connections, string theory comparison, swampland distance conjecture, species scale, tower/lattice weak gravity conjecture, de Sitter swampland, asymptotic darkness, black hole entropy, information paradox, island formula, quantum extremal surface, holographic complexity, tensor network, MERA, error correction, fault tolerance, topological quantum computing, anyon models, modular tensor categories, quantum simulation, analogue gravity, Unruh-DeWitt detector, dynamical Casimir, quantum friction, photon-photon scattering, light-by-light, Euler-Heisenberg, Born-Infeld, noncommutative QFT, fuzzy sphere, Moyal product, Seiberg-Witten map, spectral triple classification, real spectral triples, KO-dimension, Poincar\'{e} duality, orientation, grading, first-order condition, second-order condition, higher-order condition, AC geometry, almost-commutative, product geometry, internal fluctuations, unimodularity, volume quantisation, cosmological constant from NCG, Higgs mechanism from NCG, electroweak mixing from NCG, strong interaction from NCG, flavour geometry, generation structure, family symmetry, modular flavour symmetry, finite simple groups, Monster moonshine, vertex operator algebras.


%=============================================================================
\section{Publication Pipeline Update}
%=============================================================================

\begin{center}
\begin{tabular}{clcc}
\toprule
\textbf{\#} & \textbf{Paper} & \textbf{Status} & \textbf{Target} \\
\midrule
1 & $C_\ell$ (DFD vs Planck) & \textbf{SUBMITTED (arXiv)} & June 2026 \\
2 & Euclid pre-registration & 100\% & July 2026 \\
3 & No-screening & 100\% & Late July 2026 \\
4 & Unified $\alpha$ & 100\% & Aug--Sep 2026 \\
5 & \textit{Phys.\ Reports} review & 100\% & Sep 2026 \\
6 & PRL letter & \textbf{100\% (NEW at i71)} & Sep--Oct 2026 \\
7 & DESI JCAP & 100\% & August 2026 \\
\midrule
S1 & Software/JOSS & \textbf{100\% (NEW at i71)} & After papers \\
\midrule
8 & $E_G$ standalone & 55\% & Oct--Nov 2026 \\
9 & CKM paper & Outline (from v3.3 Sec 15) & Q1 2027 \\
10 & $g-2$ paper & Outline (from v3.3 Sec 16) & Q1 2027 \\
11 & v3.3 flagship update & 80\% (16/20) & October 2026 \\
\bottomrule
\end{tabular}
\end{center}

\end{document}
